\begin{proof}
Suppose $\fleetAssignVar$ is feasible for the relaxation $\model{\discretization}$ but infeasible for the original problem. The constraints of the original problem are:
1. Flow balance and demand satisfaction (same structure as LBM).
2. Resource availability (ground inventory $\ge 0$).
3. Service separation (inter-departure time $\ge \minInterdepTime_{ij}$).

If the infeasibility is due to service separation, then there exists a segment $(i,j)$ and two services departing at $t_1, t_2$ such that $|t_1 - t_2| < \minInterdepTime_{ij}$. In the LBM, these services are assigned to arcs in $\net$. If they are assigned to the same time point or close time points such that their interval overlaps in the original time scale, the calculation of $\sepViolFlow_{ijt}$ over the interval $[t, t+\minInterdepTime_{ij})$ will capture this. Specifically, since the LBM allows them (satisfying relaxed constraints), but they violate the original constraint, they must fall into a window of size $\minInterdepTime_{ij}$ with a count $> 1$. Thus, a separation violation is detected.

If the infeasibility is due to resource availability, it means that at some node $j$ and time $t$, a service is scheduled to depart but no resource unit is available. Since $\fleetAssignVar$ satisfies flow balance in $\net$, the resource unit must have "arrived" in the model. However, in reality, it has not. This implies that the model arrival time $t'_{arr}$ satisfies $t'_{arr} \le t$, but the actual arrival time $T_{arr}$ satisfies $T_{arr} > t$. This is exactly the condition $\cumDepart_{jf}(t) > \cumArrReal_{jf}(t) + \fleetCountStartByI_{jf}$. Thus, an impossible departure violation is detected.

Since the infeasibility must stem from one of the relaxed constraints (time aggregation or travel time shortening), one of these two detection mechanisms will be triggered.
\end{proof}